\subsection{Synthetic Function: Experimental Setup}
\label{sec:synthetic-setup}

\paragraph{Task family and split.}
We study a two-dimensional, task-parameterized Branin function.  For
$x=(x_1,x_2)\in[-5,10]\times[0,15]$ and task parameter $t\in[0,1]$, the
minimization objective is
\begin{equation}
 f_t(x)=\left(x_2-\frac{5.1}{4\pi^2}x_1^2+\frac{5}{\pi}x_1-6\right)^2
       +10(1-t)\cos(x_1)+10.
\end{equation}
The optimizer receives the score $y_t(x)=-f_t(x)$, so larger scores are
better.  We use 50 training tasks with
$t_i=0.005+0.02i$ for $i=0,\ldots,49$ and 20 held-out tasks with
$t_j=0.05j$ for $j=0,\ldots,19$.  Thus the held-out task parameters
interleave with, rather than duplicate, the training parameters.  The
optimum is $f_t^\star=10t$, which permits exact simple-regret evaluation.

\paragraph{Sequential optimization and model training.}
At each held-out task, an arm starts with five initial points and then
makes 50 additional, sequential oracle calls.  The standard BO loop fits
an RBF-kernel GP to the observations and maximizes a GP-UCB acquisition
function over 128 uniformly sampled candidate points per call.  Coordinates
are scaled to the unit box before GP fitting; the kernel lengthscale is
$0.2$ and the UCB exploration coefficient is $2.0$.  The BOLT and ORPT
proposal models are initialized from Qwen2.5-3B-Instruct.  At training
milestones $m\in\{2,5,10,20,30,40,50\}$, BOLT fine-tunes on the 20
highest-scoring observed points from each of the first $m$ training tasks,
using five SFT epochs.  ORPT applies one DPO epoch to the corresponding
BOLT checkpoint.  BOLT and ORPT use separate training-trajectory chains;
their held-out evaluations use the same task parameters and BO budget.

\paragraph{Matched-intervention preferences.}
For each training task at a milestone, we reserve at most 30 candidate
initial points from its observed trajectory.  We sample three shared
background pools of four points each from the remaining observed points,
with probabilities proportional to the reference model likelihood at
temperature $1$.  Each candidate is inserted into each background and
followed by $H$ BO acquisition steps. The figures reported here use
the completed $H=1$ run; separate $H=2$ and $H=3$ configurations test
longer horizons.  The candidate utility is the best score among the
inserted candidate and the points acquired in those $H$ steps; the shared
background points condition the BO trajectory but do not enter this
utility.  Matched candidate differences are evaluated with the same
backgrounds and random seeds.  We retain up to 30 preference pairs per
task with an absolute paired $z$-score of at least $1$, requiring the
preferred candidate's own observed score to exceed the rejected
candidate's score.  Preference construction uses a separate GP acquisition
configuration: 512 uniformly sampled candidates, lengthscale $0.05$, and
UCB coefficient $0$.

\paragraph{Baselines.}
We compare against STBO, POGPE, OptFormer, and LLAMBO.
